\documentclass[11pt]{article}
\usepackage[letterpaper,margin=0.8in]{geometry}
\usepackage[T1]{fontenc}
\usepackage{mathptmx}
\usepackage{tabularx,booktabs,array}
\usepackage{enumitem}
\usepackage[table]{xcolor}
\usepackage{titlesec}
\usepackage[hidelinks]{hyperref}
\usepackage{fancyhdr}
\usepackage{needspace}
\definecolor{cmured}{HTML}{C41230}
\definecolor{rowgray}{HTML}{F3F3F3}
\titleformat{\section}{\large\bfseries\color{cmured}}{\thesection.}{0.5em}{}
\titleformat{\subsection}{\normalsize\bfseries}{\thesubsection}{0.5em}{}
\titlespacing*{\section}{0pt}{10pt}{3pt}
\titlespacing*{\subsection}{0pt}{6pt}{2pt}
\setlist{itemsep=1pt,topsep=2pt,leftmargin=1.3em}
\setlength{\parindent}{0pt}
\setlength{\parskip}{4pt}
\renewcommand{\arraystretch}{1.25}
\newcolumntype{L}{>{\raggedright\arraybackslash}X}
\newcolumntype{P}[1]{>{\raggedright\arraybackslash}p{#1}}
\pagestyle{fancy}
\fancyhf{}
\rhead{\small Group 19 \textbullet{} Intro to AI, Fall 2026}
\cfoot{\small\thepage}
\renewcommand{\headrulewidth}{0.4pt}

\pagestyle{empty}
\begin{document}

\begin{center}
{\LARGE\bfseries Bringing the Collaborative Discussion Toolkit Home}\\[4pt]
{\large How a Carnegie Mellon student team is helping move 46 activities to the IF Website, word for word}
\end{center}

\vspace{4pt}
The Collaborative Discussion Toolkit is a free collection of 46 activities that help groups listen, explore ideas, and work through questions together. The activities are organized in five modules: Introduction, Creative, Critical, Culturally Responsive, and Civic Collaboration. Until now, the Toolkit has lived on its own separate website. This fall, the Interactivity Foundation decided to bring it home to the IF Website, so that all of IF's work is in one place and the Toolkit is easier to keep up to date.

\section*{The challenge}
Moving 46 activities by hand takes time, and small mistakes matter. A changed step time, a missing journal prompt, or one added word in a facilitation step can change how an activity runs. IF is also revising activities and adding new ones, so the Toolkit needs to stay easy to update after the move.

\section*{What the team built}
A team of five students from Carnegie Mellon University's Heinz College built a tool that copies each activity exactly, with its PDF and images, into one consistent page layout on the IF Website. The wording is copied by the program, not rewritten by AI. Every word, step, time, and link is then checked against the original. If anything differs, that activity is set aside and the difference is shown to a person.

Every activity arrives as a draft. Nothing appears on the website until IF staff have reviewed it and chosen to publish it. AI helpers coordinate the steps and explain in plain words why an activity was set aside, but they cannot publish anything.

\section*{Keeping the Toolkit growing}
IF staff will also have a simple, guided way to add new activities and revise existing ones. They can use it with an AI assistant or entirely by hand, with a fill-in template and a review checklist that keep every activity in the same layout.

\section*{Tested with care}
Before touching the real Toolkit, the team tested the tool on a practice set of 46 invented activities built to match the real Toolkit's structure. All 46 were copied exactly. The team then planted nine kinds of errors, including a changed step time, a dropped step, a reworded learning goal, a single added ``not,'' and a stray email address. With all of the safety checks switched on, none of them reached a draft.

\section*{What comes next}
IF staff will review the activities and publish them on the IF Website. Once they are published, links to the old Toolkit pages can lead visitors to the new ones.

\section*{About the team}
Joshua Kalakoti, Like Cui, Nibedita Biswal, Letu Wang, and Yinfeng Liu worked with the Interactivity Foundation as part of the course Intro to AI at Carnegie Mellon University's Heinz College, Fall 2026.

\end{document}
